\documentclass[11pt,reqno]{amsart}
\usepackage[margin=1in]{geometry}
\usepackage{amsmath,amssymb,amsthm}
\usepackage[colorlinks=true,linkcolor=blue,citecolor=blue,urlcolor=blue]{hyperref}

\theoremstyle{plain}
\newtheorem{proposition}{Proposition}
\newtheorem{lemma}[proposition]{Lemma}
\newtheorem{claim}[proposition]{Claim}
\newtheorem{corollary}[proposition]{Corollary}
\theoremstyle{remark}
\newtheorem{remark}[proposition]{Remark}

\newcommand{\dd}{\,\mathrm{d}}
\newcommand{\R}{\mathbb{R}}
\newcommand{\Prob}{\mathbb{P}}
\newcommand{\eps}{\varepsilon}
\newcommand{\pinst}{p_\star}
\newcommand{\OK}{{\footnotesize\textsf{[computed here]}}}
\newcommand{\cited}{{\footnotesize\textsf{[classical, cited]}}}
\newcommand{\gap}{{\footnotesize\textsf{[remaining]}}}

\begin{document}

\title[Step~1: the isomonodromy identification]{Step~1 toward O1$^\star$: the fluctuation
operator carries the Painlev\'e~II signature, and its $a$-flow is isomonodromic}
\author{Solomon Williams}
\date{\today}

\begin{abstract}
Continuation of the O1 note. O1$^\star$ asks that the weak-noise fluctuation coefficients
$c_n(a)$ of the $\mathrm{TW}_\beta$ tail be the trans-series coefficients of Painlev\'e~II, so
that Costin's Borel-summability theorem applies. This note carries out the second variation of the
Freidlin--Wentzell action about the Airy instanton and finds the fluctuation (Jacobi) operator
\[
\mathcal J=-\partial_x^2+\bigl(6\pinst^2-2(x+a)\bigr),
\]
whose $6\pinst^2$ term is exactly the coefficient of the Painlev\'e~II linearization
$-\partial_s^2+(6q^2+s)$ --- a nontrivial fingerprint, not a coincidence of normalization. It then
frames the identification correctly: the fixed-level linear problem is Airy (trivial deformation),
and the Painlev\'e~II structure lives in the \emph{spectral-level} variable $a$, where the
Tracy--Widom theorem makes the $\beta=2$ tail \emph{exactly} the Hastings--McLeod
$\tau$-function. Step~1 thereby reduces to one precise obligation --- identify the fluctuation
determinant / transport hierarchy with the Jimbo--Miwa $\tau$-function of the $a$-isomonodromic
Painlev\'e~II family --- for which the machinery (Flaschka--Newell linear problem, JMU deformation
theory) is classical and Costin then closes summability.
\end{abstract}

\maketitle

\section{Setup and the goal of Step~1}

From the O1 note: the tail $\mathcal T_\beta(a)=\Prob(\text{Riccati explosion})$ has the
weak-noise (large-$\beta$, $\eps^2=1/\beta$) form
$\mathcal T_\beta(a)=e^{-\beta\Phi(a)}\sum_{n\ge0}c_n(a)\,\beta^{-n}$, where the eikonal instanton
$\pinst(x)$ solves the fluctuation-branch Riccati equation
\begin{equation}\label{eq:inst}
\dot\pinst=\pinst^2-(x+a),\qquad \pinst=-\varphi'/\varphi,\quad \varphi''=(x+a)\varphi\ \ (\text{Airy}),
\end{equation}
and $\Phi(a)=\tfrac12\int b^2\dd x$ with $b:=(x+a)-\pinst^2=-\dot\pinst$. The $c_n$ are generated by
the transport hierarchy of the fluctuation operator $\mathcal J:=\tfrac12\delta^2 I[\pinst]$, the
Hessian of the Freidlin--Wentzell action $I[p]=\tfrac18\int(\dot p-b)^2\dd x$. Step~1 is to prove:
\begin{quote}\itshape
$\mathcal J$ is the linearization of Painlev\'e~II about its Hastings--McLeod solution, in the
spectral variable $a$; equivalently, the $\{c_n(a)\}$ are the PII trans-series coefficients. Then
Costin's theorem gives Borel summability with singularities at the instanton action.
\end{quote}

\section{The fluctuation operator (explicit)}\label{sec:jacobi}

\begin{proposition}[Jacobi operator]\label{prop:jacobi}
Write $F[p]:=\dot p-b=\dot p-(x+a)+p^2$, so $I=\tfrac18\int F^2$. About the instanton
\eqref{eq:inst}, with $\eta:=\delta p$,
\[
\delta^2 I[\pinst]=\tfrac14\!\int \eta\,\mathcal J\,\eta\,\dd x,\qquad
\boxed{\ \mathcal J=-\partial_x^2+Q(x),\quad Q(x)=6\pinst^2-2(x+a).\ }
\]
\end{proposition}

\begin{proof}
$\delta F=\dot\eta+2\pinst\eta$ and $\delta^2F=2\eta^2$, so
$\delta^2 I=\tfrac18\int\!\bigl[2(\delta F)^2+2F_\star\,\delta^2F\bigr]
=\tfrac14\!\int(\dot\eta+2\pinst\eta)^2+\tfrac12\!\int F_\star\eta^2$,
where on the instanton $F_\star=\dot\pinst-b_\star=-2b_\star$ (since $\dot\pinst=-b_\star$).
Integrating $\int(\dot\eta+2\pinst\eta)^2$ by parts,
$\int\dot\eta^2=-\int\eta\ddot\eta$ and $\int4\pinst\eta\dot\eta=-2\int\dot\pinst\eta^2$, gives
$\tfrac14\int\eta\bigl[-\ddot\eta+(4\pinst^2-2\dot\pinst)\eta\bigr]-\int b_\star\eta^2$. Using
$\dot\pinst=\pinst^2-(x+a)$ and $b_\star=(x+a)-\pinst^2$,
$4\pinst^2-2\dot\pinst-4b_\star=6\pinst^2-2(x+a)$. \qedhere
\end{proof}

\begin{corollary}[The Painlev\'e~II fingerprint]\label{cor:fingerprint}
The Painlev\'e~II equation $q_{ss}=2q^3+sq$ linearizes about a solution $q$ as
$\delta_{ss}=(6q^2+s)\delta$, i.e.\ the operator $-\partial_s^2+(6q^2+s)$. The fluctuation operator
of Proposition~\ref{prop:jacobi} has \emph{exactly} the PII coefficient $6(\cdot)^2$ on its
field-squared term. A generic weak-noise Hessian would carry an arbitrary constant there; the value
$6$ is forced and is the isomonodromy fingerprint. \OK
\end{corollary}

\begin{remark}[What the two remaining discrepancies mean]\label{rem:disc}
$\mathcal J=-\partial_x^2+6\pinst^2-2(x+a)$ differs from $-\partial_s^2+6q^2+s$ in (i) the sign/scale
of the linear-in-position term and (ii) the field: $\pinst$ solves \emph{Airy}, not PII. Both are
resolved by the correct framing of \S\ref{sec:iso}: the linear term and the field belong to the
\emph{fixed-level} ($x$) problem, which is free (Airy); the PII lives in the deformation variable
$a$, and $\pinst$ is the $a$-family's generating field, whose flow --- not whose $x$-profile --- is
PII. Corollary~\ref{cor:fingerprint} is the trace this leaves on the $x$-Hessian.
\end{remark}

\section{The correct framing: PII lives in the spectral variable $a$}\label{sec:iso}

\begin{proposition}[Fixed level is free; the $\beta=2$ anchor is exactly PII]\label{prop:anchor}
(i) At fixed $a$ the instanton problem is the Airy equation \eqref{eq:inst}; its irregular
singularity at $x=\infty$ has $a$-independent Stokes data (translation $z=x+a$ removes $a$). Hence
the leading (Airy) deformation in $a$ is trivial, and any PII structure is carried by the
interacting fluctuations, i.e.\ by the $6\pinst^2$ term of $\mathcal J$. \cited\,/\,\OK
(ii) At $\beta=2$ the tail is \emph{exactly} Painlev\'e~II: by Tracy--Widom,
$F_2(a)=\exp\!\bigl(-\int_a^\infty (x-a)\,q(x)^2\dd x\bigr)$ with $q$ the Hastings--McLeod solution
of $q_{ss}=2q^3+sq$, and $\mathcal T_2(a)=1-F_2(a)$. \cited
\end{proposition}

Thus $a$ is the Painlev\'e (deformation) variable and $x$ is the internal/spectral variable of the
associated linear problem. The classical structure that governs this is:
\begin{itemize}
\item[(FN)] the Flaschka--Newell $2\times2$ linear system whose isomonodromic deformation is PII,
with an irregular singularity carrying the Stokes/connection data; \cited
\item[(JMU)] the Jimbo--Miwa--Ueno $\tau$-function of that isomonodromic family, whose logarithmic
derivative in the deformation variable is the Hamiltonian of PII. \cited
\end{itemize}
The Airy problem \eqref{eq:inst} is the degeneration of the FN linear system in which the PII field
is switched off; turning on the fluctuations ($6\pinst^2$) deforms it, and the deformation in $a$
is the PII flow.

\section{Reduction of Step~1 to one $\tau$-function identity}

\begin{claim}[Step~1 reduced]\label{cl:reduce}
Let $\mathcal D(a;\beta):=\det\nolimits_\zeta\mathcal J$ be the (zero-mode-removed, O4)
fluctuation determinant of $\mathcal J$, and let $\tau_{\mathrm{JMU}}(a)$ be the Jimbo--Miwa
$\tau$-function of the $a$-isomonodromic PII family of \S\ref{sec:iso}. If
\begin{equation}\label{eq:tauid}
\sum_{n\ge0}c_n(a)\,\beta^{-n}\ \dot=\ \text{the large-$\beta$ asymptotic expansion of }
\tau_{\mathrm{JMU}}(a)\ \text{along the Hastings--McLeod trans-series,}
\end{equation}
then the $c_n(a)$ are the PII trans-series coefficients, and Costin's theorem
(rank-one nonlinear ODE at an irregular singularity) gives their Borel summability in $1/\beta$
with singularities at the instanton action, closing O1$^\star$ and hence Theorem~T1. \gap
\end{claim}

Two independent supports that \eqref{eq:tauid} is the right target, beyond
Corollary~\ref{cor:fingerprint}:
\begin{enumerate}
\item[(S1)] \emph{Exact at $\beta=2$.} Proposition~\ref{prop:anchor}(ii): there
$\mathcal T_2(a)$ \emph{is} the Hastings--McLeod expression, so \eqref{eq:tauid} holds identically
(not merely asymptotically) at $\beta=2$, fixing the $\tau$-function normalization the large-$\beta$
expansion must match. \cited
\item[(S2)] \emph{Correct reality/Stokes structure.} The Hastings--McLeod trans-series has
factorially divergent, non-sign-alternating coefficients in the relevant region, with the
median-resummation reality condition (Cleri--Dunne); a linear transport hierarchy has no reason to
reproduce that nonlinear Stokes structure, whereas a $\tau$-function identification does. \cited
\end{enumerate}

\paragraph{A concrete sub-computation toward \eqref{eq:tauid}.}
The $a$-derivative of the instanton, $\psi:=\partial_a\pinst$, is a Jacobi field: differentiating
\eqref{eq:inst} gives the linear inhomogeneous equation
\begin{equation}\label{eq:jacobifield}
\dot\psi=2\pinst\psi-1.
\end{equation}
Its homogeneous part is annihilated by the same operator whose second-order form is
$\mathcal J$; the inhomogeneity $-1$ is the source that, iterated through the transport hierarchy,
generates the $a$-dependence of the $c_n$. The precise statement to establish is that the pair
$(\pinst,\psi)$ solving \eqref{eq:inst}--\eqref{eq:jacobifield} maps to the Flaschka--Newell frame
of PII, so that $\partial_a\log\mathcal D=\partial_a\log\tau_{\mathrm{JMU}}$ up to the explicit
one-loop term. This is a finite (Lax-pair) identification, the natural next write-up. \gap

\section{Status}

\emph{Computed here:} the fluctuation operator $\mathcal J=-\partial_x^2+6\pinst^2-2(x+a)$
(Prop.~\ref{prop:jacobi}) and its Painlev\'e~II fingerprint (Cor.~\ref{cor:fingerprint});
the Jacobi-field equation \eqref{eq:jacobifield}.
\emph{Classical, cited (the backbone):} Airy triviality at fixed level and the Tracy--Widom
$\beta=2$ PII identity (Prop.~\ref{prop:anchor}); the Flaschka--Newell linear problem and the
JMU $\tau$-function.
\emph{Remaining (the crux of Step~1):} the $\tau$-function identity \eqref{eq:tauid}, reduced to the
finite Lax-pair identification of the sub-computation. Once it holds, Costin closes summability.
\emph{Assessment:} the fingerprint (Cor.~\ref{cor:fingerprint}) plus the exact $\beta=2$ anchor
(S1) make the PII identification very likely correct; Step~1 is now a bounded isomonodromy
computation, not an open-ended search. Recommended next write-up: the Lax-pair map of
\eqref{eq:inst}--\eqref{eq:jacobifield} into the Flaschka--Newell frame.

\begin{thebibliography}{9}
\bibitem{TW1994} C.~A.~Tracy, H.~Widom, \emph{Level-spacing distributions and the Airy kernel},
Comm. Math. Phys. \textbf{159} (1994) 151--174.
\bibitem{FN1980} H.~Flaschka, A.~C.~Newell, \emph{Monodromy- and spectrum-preserving deformations
I}, Comm. Math. Phys. \textbf{76} (1980) 65--116.
\bibitem{JMU1981} M.~Jimbo, T.~Miwa, K.~Ueno, \emph{Monodromy preserving deformation of linear
ordinary differential equations with rational coefficients I}, Physica D \textbf{2} (1981)
306--352.
\bibitem{CleriDunne} N.~J.~Cleri, G.~V.~Dunne, \emph{Resurgent trans-series for generalized
Hastings--McLeod solutions}, J. Phys. A \textbf{53} (2020) 305201; arXiv:2002.06270.
\bibitem{CLDS} A.~Comtet, P.~Le Doussal, N.~R.~Smith, \emph{The Tracy--Widom distribution at large
Dyson index}, arXiv:2510.14433 (2025).
\bibitem{Costin} O.~Costin, \emph{On Borel summation and Stokes phenomena of nonlinear
differential systems}, Duke Math. J. \textbf{93} (1998) 289--344; arXiv:math/0608408.
\end{thebibliography}

\end{document}
